% =====================================================================
% Complexity analysis and factorization bound for Hierarchical Block-NCM
% Every number in this section is reproduced by
%   py paper/scripts/complexity_constants.py [--N --M --C --K --d --Lx]
% =====================================================================

\section{Complexity Analysis and Factorization Bound}
\label{sec:complexity}

\subsection{Setup and notation}
\label{sec:notation}

A prompt $X$ of length $L_x$ is followed by an output $Y$ of $N$ tokens.
Generation proceeds over $K$ hierarchy levels $k\in\{1,\dots,K\}$ with scale
factors $s_k\in(0,1]$, $s_K=1$: level $k$ of the monolithic model contains
$s_kN$ positions. We write
\begin{equation}
  S_k=\sum_{j\le k}s_j,\qquad
  \sigma_1=\sum_{k=1}^{K}s_k,\qquad
  \sigma_2=\begin{cases}
    \sum_{k}s_kS_k=\tfrac12\big(\sigma_1^2+\sum_{k}s_k^2\big) & \text{block-causal},\\[2pt]
    \sum_{k}s_k^2 & \text{block-diagonal}.
  \end{cases}
  \label{eq:sigmas}
\end{equation}
Here $\sigma_2N^2$ counts the query--key pairs between output positions
(Assumption~\ref{as:mask}). The block-causal identity holds because
$\sum_k s_kS_k$ sums over all pairs $j\le k$, so it equals half the sum over
all ordered pairs plus half the diagonal.
For the geometric schedule $s_k=2^{-(K-k)}$ with block-causal attention,
$\sigma_1=2-2^{1-K}$ and $\sigma_2=\tfrac12\big[\sigma_1^2+\tfrac43(1-4^{-K})\big]$; at
$K=6$, $\sigma_1=1.969$, $\sigma_2=2.604$ and $\sigma_1/\sigma_2=0.756$.
The released NCM TinyStories configuration generates levels of length
$(8,16,24,32,40,52,64,76,92,108,128)$ after a single-position level~0, with
block-diagonal attention; there $N=128$ latent positions, $\sigma_1=5.000$ and
$\sigma_2=3.211$.

\emph{Block-NCM} splits $Y$ into $M$ blocks $B_1,\dots,B_M$ of equal length
$N/M$. Block $m$ is conditioned on $X$ and on a global blueprint
$P=(P_1,\dots,P_M)$ of $MC$ discrete codes ($C$ per block). All blocks are
folded into the batch dimension and decoded concurrently.

\paragraph{Cost model.} A transformer layer of width $d$ spends
$c_{\mathrm{lin}}d^2$ FLOPs per processed token on projections and MLP
($c_{\mathrm{lin}}=24$ for a $4\times$ MLP, counting 2 FLOPs per multiply-add) and
$4d$ FLOPs per query--key pair on attention ($QK^\top$ and $AV$). We ignore
softmax, normalization and embedding costs. All quantities below are per layer;
they scale linearly with depth, which cancels in every ratio.

\paragraph{Assumptions.}
\begin{enumerate}[label=(A\arabic*),leftmargin=*,itemsep=1pt]
  \item \label{as:mask} \emph{Level-wise generation.} At level $k$, only the
  $s_kN$ positions of that level issue queries. They attend to the prompt and
  either to all positions of levels $j\le k$, whose keys and values are cached
  (\emph{block-causal}), or to level $k$ only (\emph{block-diagonal}). The
  released NCM configurations use block-diagonal attention and no cross-level
  cache: earlier levels reach level $k$ only through the running reconstruction
  $\hat z$, area-pooled to length $s_kN$ (\texttt{models/next\_concept.py},
  \texttt{models/quant.py}).
  \item \label{as:block} \emph{Block-local attention.} In Block-NCM, the level-$k$
  queries of block $m$ attend to $X$, to the full blueprint $P$, and to the
  same set of levels as in \ref{as:mask}, restricted to block $m$.
  \item \label{as:equal} \emph{Equal block lengths.} All blocks have length $N/M$;
  variable-length blocks would need padding or length codes in $P$.
\end{enumerate}

\subsection{Attention FLOPs}

\begin{proposition}[Attention cost]
\label{prop:attn}
Under \ref{as:mask}--\ref{as:equal} with block-causal attention, the per-layer
attention FLOPs are given below. For block-diagonal attention, replace $S_k$ by
$s_k$; the closed forms are unchanged with $\sigma_2$ taken from
\eqref{eq:sigmas}.
\begin{align}
  A_{\mathrm{mono}}
  &= 4d\sum_{k=1}^{K}(s_kN)\big(L_x+S_kN\big)
   = 4d\big(N^2\sigma_2+NL_x\sigma_1\big),
  \label{eq:attn-mono}\\
  A_{\mathrm{block}}
  &= 4d\,M\sum_{k=1}^{K}\Big(s_k\frac{N}{M}\Big)\Big(L_x+MC+S_k\frac{N}{M}\Big)
   = 4d\Big(\frac{N^2\sigma_2}{M}+N\,(L_x+MC)\,\sigma_1\Big).
  \label{eq:attn-block}
\end{align}
\end{proposition}

\begin{proof}
At level $k$ the monolithic model has $s_kN$ queries, each attending to $L_x$
prompt keys and $\sum_{j\le k}s_jN=S_kN$ output keys, at $4d$ FLOPs per pair;
summing over $k$ gives \eqref{eq:attn-mono}. In Block-NCM each of the $M$
blocks has $s_kN/M$ queries at level $k$, attending to $L_x+MC$ shared keys and
$S_kN/M$ block-local keys. Multiplying by $M$ and summing over $k$ gives
\eqref{eq:attn-block}.
\end{proof}

\begin{corollary}[Attention ratio and optimal block count]
\label{cor:ratio}
With $L_x=0$,
\begin{equation}
  r_{\mathrm{attn}}
  \;=\;\frac{A_{\mathrm{block}}}{A_{\mathrm{mono}}}
  \;=\;\frac1M+\frac{MC}{N}\cdot\frac{\sigma_1}{\sigma_2}.
  \label{eq:ratio}
\end{equation}
For fixed $K$, $\{s_k\}$ and $C$, $r_{\mathrm{attn}}\to 1/M$ if and only if
$M^2C=o(N)$. The block count that minimizes $A_{\mathrm{block}}$ is
\begin{equation}
  M^\star=\sqrt{\frac{N\sigma_2}{C\sigma_1}},
  \label{eq:mstar}
\end{equation}
and it does not depend on $L_x$.
\end{corollary}

\begin{proof}
Divide \eqref{eq:attn-block} by \eqref{eq:attn-mono} with $L_x=0$. For
\eqref{eq:mstar}, the part of $A_{\mathrm{block}}$ that depends on $M$ is
$N^2\sigma_2/M+NMC\sigma_1$, which is convex in $M>0$; setting its derivative
to zero gives $M^\star$.
\end{proof}

For $N=2048$, $M=8$, $C=8$, $K=6$ (geometric), \eqref{eq:ratio} gives
$r_{\mathrm{attn}}=0.125+0.024=0.149$: a $6.7\times$ reduction rather than
$8\times$, with $M^\star\approx18$. A non-zero prompt weakens the reduction:
$L_x=256$ gives $r_{\mathrm{attn}}=0.222$ ($4.5\times$).

\subsection{End-to-end FLOPs}

Attention is only part of the cost. The monolithic model processes $N\sigma_1$
output tokens, plus $L_x$ prompt tokens during prefill. Block-NCM processes the
same output tokens plus the $MC$ blueprint codes (and $ML_x$ prompt tokens
unless the prompt prefix is shared across the batch). Let
\begin{equation}
  \phi=\frac{A_{\mathrm{mono}}}{A_{\mathrm{mono}}+c_{\mathrm{lin}}d^2N\sigma_1}
  =\frac{N\sigma_2}{N\sigma_2+6d\sigma_1}\qquad(L_x=0,\ c_{\mathrm{lin}}=24)
  \label{eq:phi}
\end{equation}
be the attention share of monolithic FLOPs. Then
\begin{equation}
  r_{\mathrm{total}}
  =(1-\phi)\Big(1+\frac{MC}{N\sigma_1}\Big)+\phi\,r_{\mathrm{attn}}.
  \label{eq:rtotal}
\end{equation}
Because $\phi<1$ whenever $N$ is not much larger than $d$, the end-to-end
reduction is far smaller than $M$ (Table~\ref{tab:constants}). The stage-0
planner and the boundary verifier add further cost on top of
\eqref{eq:rtotal}; both must be measured and reported.

\subsection{KV-cache and activation memory}

Let $b$ be the bytes per element and $g\le1$ the grouped-query reduction
factor. The KV cache at the last level holds $2gdb$ bytes per resident token
per layer.

\begin{proposition}[KV cache is not reduced]
\label{prop:kv}
When all $M$ blocks are resident concurrently, the number of cached tokens is
\begin{align}
  T_{\mathrm{mono}}&=L_x+N_{\mathrm{res}},\\
  T_{\mathrm{block}}&=\underbrace{L_x\ \text{or}\ ML_x}_{\text{prefix shared / copied}}
  +\underbrace{MC\ \text{or}\ M^2C}_{\text{blueprint shared / copied}}
  +N_{\mathrm{res}}\;\;\ge\;T_{\mathrm{mono}},
\end{align}
where $N_{\mathrm{res}}=N\sigma_1$ for block-causal attention with a cross-level
cache and $N_{\mathrm{res}}=N$ for block-diagonal attention.
\end{proposition}

\begin{proof}
Each output position at each level is cached exactly once in both models,
because block-local attention changes \emph{which} keys a query reads, not how
many keys exist. Block-NCM additionally caches the blueprint, and also copies
the prompt and the blueprint per block unless they are shared.
\end{proof}

\begin{remark}[Where the $1/M$ appears in memory]
The quantity that does shrink by roughly $M$ is the attention score matrix at
level $K$: $N(L_x+N\sigma_1)$ entries per head for the monolithic model, versus
$N(L_x+MC)+N^2\sigma_1/M$ for Block-NCM. Fused kernels such as
FlashAttention or PyTorch SDPA never materialize this matrix, so with them
peak memory is dominated by weights, activations and the KV cache, and we
expect no material reduction in peak VRAM from blocking.
Experiment~1 therefore reports peak memory with fused and with unfused
attention separately.
\end{remark}

\begin{remark}[Sequential passes]
Monolithic NCM performs $K$ sequential forward passes. Block-NCM performs $K$
passes plus the stage-0 planner's passes. If the planner emits the $MC$
blueprint codes autoregressively, it adds up to $MC$ sequential passes, which
exceeds $K$ for typical settings ($MC=64$). A single-pass or $O(\log(MC))$-pass
planner is required for the latency argument to hold.
\end{remark}

\begin{table}[t]
\centering
\caption{Analytical cost ratios, Block-NCM vs.\ monolithic NCM ($M=8$).
Top: hypothetical long-sequence model ($C=8$, $K=6$, geometric scales,
block-causal). Bottom: the released TinyStories configuration ($d=768$,
block-diagonal, levels $8,\dots,128$, $L_x=1$, $C=1$). Ratios below 1 favour
Block-NCM. Reproduced by \texttt{scripts/complexity\_constants.py}.}
\label{tab:constants}
\small
\begin{tabular}{@{}rrrcccc@{}}
\toprule
$N$ & $d$ & $L_x$ & $r_{\mathrm{attn}}$ & $\phi$ & $r_{\mathrm{total}}$ & KV ratio (shared / copied)\\
\midrule
2048 & 1024 & 0   & 0.149 & 0.306 & 0.751 & 1.016 / 1.127\\
2048 &  768 & 0   & 0.149 & 0.370 & 0.695 & 1.016 / 1.127\\
2048 & 1024 & 256 & 0.222 & 0.312 & 0.767 & 1.015 / 1.537\\
8192 &  768 & 0   & 0.131 & 0.702 & 0.391 & 1.004 / 1.032\\
\midrule
 128 &  768 & 1   & 0.232 & 0.018 & 0.999 & 1.062 / 1.550\\
\bottomrule
\end{tabular}
\end{table}

\begin{remark}[Released NCM checkpoints]
At the released scale ($N=128$ latent positions, $d=768$), attention is
$\phi\approx1.8\%$ of generator FLOPs. Block partitioning still cuts attention
by $4.3\times$, but saves under $1\%$ of FLOPs end to end for any
$M\in\{2,4,8\}$. The complexity benefit of
Block-NCM therefore appears only for outputs much longer than the
128-position context of the released tokenizers.
\end{remark}

% ---------------------------------------------------------------------
\subsection{Quality cost of block factorization}
\label{sec:kl}

Let $p$ be the data distribution and assume the blueprint is a deterministic
encoding of the output, $P=f(Y)$ (e.g.\ the stage-0 tokenizer applied to the
ground-truth blocks). The model is
\begin{equation}
  q(P,Y\mid X)=q_\theta(P\mid X)\prod_{m=1}^{M}q_\theta(B_m\mid X,P),
  \qquad q(Y\mid X)=\sum_{P}q(P,Y\mid X).
\end{equation}
Write $p_m(\cdot\mid X,P)$ for the true marginal of block $m$ given $(X,P)$.

\begin{theorem}[Three-term decomposition]
\label{thm:kl}
\begin{align}
  \mathbb{E}_X\,D_{\mathrm{KL}}\big(p(Y\!\mid\!X)\,\Vert\,q(Y\!\mid\!X)\big)
  \;\le\;&
  \underbrace{\mathbb{E}_X\,D_{\mathrm{KL}}\big(p(P\!\mid\!X)\,\Vert\,q_\theta(P\!\mid\!X)\big)}_{\text{planner error}}
  \;+\;
  \underbrace{\mathbb{E}_{X,P}\sum_{m=2}^{M}I\big(B_m;B_{1:m-1}\mid X,P\big)}_{\text{factorization cost (total correlation)}}
  \nonumber\\
  &+\;
  \underbrace{\mathbb{E}_{X,P}\sum_{m=1}^{M}D_{\mathrm{KL}}\big(p_m(\cdot\!\mid\!X,P)\,\Vert\,q_\theta(B_m\!\mid\!X,P)\big)}_{\text{block-model error}}.
  \label{eq:kl}
\end{align}
\end{theorem}

\begin{proof}
Since $P=f(Y)$ under $p$, the joint $p(P,Y\mid X)=p(Y\mid X)\,\mathbb{1}[P=f(Y)]$
has $Y$-marginal $p(Y\mid X)$. Marginalizing out $P$ cannot increase KL
divergence (data-processing inequality), so
$D_{\mathrm{KL}}(p(Y|X)\Vert q(Y|X))\le D_{\mathrm{KL}}(p(P,Y|X)\Vert q(P,Y|X))$.
The chain rule splits the right-hand side into
$D_{\mathrm{KL}}(p(P|X)\Vert q_\theta(P|X))
 +\mathbb{E}_{P}\,D_{\mathrm{KL}}\big(p(Y|X,P)\,\Vert\,\textstyle\prod_m q_\theta(B_m|X,P)\big)$.
For the second term, with every expectation taken under $p(\cdot\mid X,P)$,
\begin{align*}
  D_{\mathrm{KL}}\Big(p\,\Big\Vert\,\prod_m q_\theta\Big)
  &=\mathbb{E}\Big[\log\frac{p(B_{1:M})}{\prod_m p_m(B_m)}\Big]
   +\sum_m\mathbb{E}\Big[\log\frac{p_m(B_m)}{q_\theta(B_m)}\Big]\\
  &=\mathrm{TC}(B_{1:M}\mid X,P)+\sum_m D_{\mathrm{KL}}(p_m\Vert q_\theta),
\end{align*}
and the chain rule of mutual information gives
$\mathrm{TC}(B_{1:M}\mid X,P)=\sum_{m=2}^{M}I(B_m;B_{1:m-1}\mid X,P)$.
Taking expectations over $X$ and $P$ completes the proof.
\end{proof}

\begin{corollary}[Irreducible cost]
\label{cor:irreducible}
Over all block decoders that factorize given $P$, the block-model error is
minimized (to zero) by $q_\theta(B_m\mid X,P)=p_m(B_m\mid X,P)$. The
total-correlation term is therefore a lower limit on what a factorized decoder
can achieve with a given blueprint, independent of model capacity.
\end{corollary}

\begin{remark}[Richer blueprints]
At the extreme $P=Y$ the total correlation is zero. In general, however,
conditioning on more information can \emph{increase} conditional mutual
information, so ``a richer blueprint reduces inter-block dependence'' is not a
theorem. It is a hypothesis about the data, and we test it empirically. The
hypothesis is most plausible for code, where function signatures and type
contracts in $P$ fix the interface between bodies, and least plausible for
prose with long-range co-reference.
\end{remark}

\paragraph{Empirical estimator.} Let $r_\psi(Y\mid X,P)$ be a
blueprint-conditioned autoregressive model trained on the same data. The
per-token negative log-likelihood (NLL) gap
\begin{equation}
  \widehat{\Delta}
  =\frac1N\,\mathbb{E}_p\Big[-\sum_m\log q_\theta(B_m\mid X,P)\Big]
  -\frac1N\,\mathbb{E}_p\big[-\log r_\psi(Y\mid X,P)\big]
\end{equation}
satisfies
$N\widehat{\Delta}=\mathrm{TC}+\sum_m D_{\mathrm{KL}}(p_m\Vert q_\theta)
 -D_{\mathrm{KL}}(p\Vert r_\psi)$, hence
\begin{equation}
  \mathbb{E}_{X,P}\,\mathrm{TC}(B_{1:M}\mid X,P)
  \;\le\;N\widehat{\Delta}+\mathbb{E}_{X,P}\,D_{\mathrm{KL}}(p\,\Vert\,r_\psi).
\end{equation}
We report $\widehat{\Delta}$ in bits per token as a function of $M$ and of the
domain (prose vs.\ code). When $r_\psi$ is a strong model, small
$\widehat{\Delta}$ is evidence that the factorization cost is small.
